% Generated draft: requires mathematical and source review.
\documentclass[11pt]{article}
\usepackage[margin=25mm]{geometry}
\usepackage{fontspec}
\setmainfont{Latin Modern Roman}
\usepackage{amsmath,amssymb,mathtools,unicode-math}
\setmathfont{Latin Modern Math}
\usepackage{enumitem,needspace,xurl,hyperref}
\setlength{\parindent}{0pt}
\setlength{\parskip}{5pt}
\setlength{\emergencystretch}{3em}
\newcommand{\N}{\mathbb{N}}
\newcommand{\Z}{\mathbb{Z}}
\newcommand{\Q}{\mathbb{Q}}
\newcommand{\R}{\mathbb{R}}
\newcommand{\C}{\mathbb{C}}
\newcommand{\F}{\mathbb{F}}
\newcommand{\A}{\mathbb{A}}
\newcommand{\PP}{\mathbb{P}}
\newcommand{\E}{\mathbb{E}}
\newcommand{\eps}{\varepsilon}
\newcommand{\dd}{\,\mathrm{d}}
\begin{document}
\section*{UnsolvedMath ID: 30002133}
\textbf{Source problem code:} catalogue entry
\subsection{Definitions and mathematical statement}
Let
\[
D=\{(x,y)\in\R^2:x^2+y^2\leq 1\},\qquad
\mu(A)=\frac{1}{\pi}\int_A\dd x\,\dd y
\]
for Lebesgue-measurable subsets \(A\subseteq D\). Consider the class
\[
\mathcal F_0=\left\{f:D\longrightarrow D:\ f\text{ is a }C^\infty\text{ diffeomorphism},\ f_*\mu=\mu,\ h_{\mathrm{top}}(f)=0\right\}.
\]
Here smoothness is understood up to the boundary, with a smooth inverse; \(h_{\mathrm{top}}(f)\) denotes topological entropy. Such a diffeomorphism maps \(\partial D\) onto itself. The supplied statement imposes no pointwise boundary-fixing condition or specified boundary dynamics. Area preservation is expressed above as preservation of standard area measure; the source record does not explicitly distinguish this convention from preservation of the oriented area form.

To spell out some possible notions of complexity, \(f\) has a dense forward orbit if
\[
\exists x\in D\quad \overline{\{f^n(x):n\geq 0\}}=D.
\]
It is ergodic with respect to \(\mu\) if every measurable set \(A\) satisfying \(\mu(f^{-1}(A)\mathbin{\triangle}A)=0\) has \(\mu(A)\in\{0,1\}\). One standard measure-theoretic definition of weak mixing is
\[
\forall A,B\subseteq D\text{ measurable},\qquad
\lim_{N\to\infty}\frac1N\sum_{n=0}^{N-1}
\left|\mu(f^{-n}(A)\cap B)-\mu(A)\mu(B)\right|=0.
\]
Strong mixing with respect to \(\mu\) means
\[
\forall A,B\subseteq D\text{ measurable},\qquad
\lim_{n\to\infty}\mu(f^{-n}(A)\cap B)=\mu(A)\mu(B).
\]
These definitions clarify possible readings of mixing; the supplied prompt does not select a particular definition of mixing or a complete list of complexity properties.

\textbf{Retained survey question.} How mixing, or otherwise dynamically complex, can a map \(f\in\mathcal F_0\) be?

\textbf{Unresolved formulation qualification.} This is a broad survey question, not a uniquely specified conjecture or existence statement. The supplied registry marks the statement as \texttt{unrecoverable}, leaves its clean statement empty, and explains that the source surveys several notions of complexity. Its verification note reports positive constructions concerning dense orbits, ergodicity, and weak mixing, but does not supply their precise hypotheses or indicate whether these claims concern one construction or several. The note assigns the source's explicit unresolved existence question to the separate record \texttt{OWR-12007-014}; that question is not substituted for the present record \texttt{catalogue entry}. No sharper mathematical target can be recovered from the supplied evidence alone.
\subsection{Short English statement}
This record asks broadly what dynamical complexity is possible for smooth area-preserving diffeomorphisms of the closed disk with zero topological entropy. The supplied verification treats it as a survey with reported positive constructions, rather than a separate precise open conjecture. Its \texttt{partially\_solved} label and literature assessment are supplied curation, not independently verified mathematical conclusions.
\subsection{Sources}
\begin{itemize}
\item Supplied local UnsolvedMath registry record, permanent ID \(30002133\), exact source code \texttt{catalogue entry}, title ``Answered Survey of Zero-Entropy Disk Dynamics,'' in the Oberwolfach Reports Open Problems collection. The project identifies \texttt{data/unsolvedmath/problems.json} and the supplementary \texttt{catalog/problems.json} as local registries; the supplied extract does not specify which file contains this record. The recorded proposed year is 2012; no individual proposer is named.
\item Supplied source citation: \emph{Geometric Group Theory, Hyperbolic Dynamics and Symplectic Geometry} (2013), Oberwolfach Reports, pp.~2139--2203, DOI 10.4171/owr/2012/35. \url{https://doi.org/10.4171/owr/2012/35}. This is also listed among the supplied literature sources as ``Oberwolfach Report 35.'' The background describes the reports as documenting workshops and their open-problem discussions.
\item Supplied literature source: ``Bramham, pseudo-rotations and zero entropy.'' \url{https://arxiv.org/abs/1205.6243}.
\item Supplied literature source: ``Bramham, periodic approximations.'' \url{https://annals.math.princeton.edu/2015/181-3/p04}.
\item Supplied literature assessment, checked exactly \texttt{2026-08-21}: status \texttt{partially\_solved}. It states that zero-entropy area-preserving disk maps can exhibit dense or weakly mixing behavior, while the strongest mixing question remains unresolved. This is supplied dated curation derived from the OWR record, not an independently verified assessment. The background expressly requests independent verification of time-sensitive status and cited claims. The supplied statement-verification note separately reports positive constructions for dense orbits, ergodicity, and weak mixing, and assigns the explicit unresolved existence question to \texttt{OWR-12007-014}.
\end{itemize}
\end{document}
